% Authored. Tables and numbers come from generated/.

\section{Results}
\label{sec:results}

\subsection{Long-run composition}

Figure~\ref{fig:longrun} plots the four balance ratios over the full window. The
visual impression is of two series that track each other closely and two that stay
near zero: General Government and Central Government are nearly indistinguishable
for long stretches, while Regional and Local Government and the Social Security
Funds move in a narrow band.

\begin{figure}[htbp]
\centering
\includegraphics[width=\linewidth]{01_long_run_balances.png}
\caption{Net lending (+) and net borrowing (--) by subsector, \PanelStart--\PanelEnd,
as a percentage of nominal GDP. The vertical rule marks the 1995 source splice; the
series either side of it come from different source families and are not chained.
Sources: Banco de Portugal / INE long series and INE via PORDATA.}
\label{fig:longrun}
\end{figure}

That closeness is itself the first result. It says the aggregate balance has been
predominantly a Central Government quantity in level terms, which is what makes the
recent divergence at the right-hand edge of the figure notable rather than routine.

Figure~\ref{fig:contributions} draws the identity directly, stacking the signed
subsector contributions against the aggregate total. Where a tall stack sits under a
shallow line, the subsectors moved in opposite directions and offset one another.

\begin{figure}[htbp]
\centering
\includegraphics[width=\linewidth]{08_subsector_contributions.png}
\caption{The identity of Equation~\eqref{eq:identity} drawn: signed subsector
contributions stacked against the General Government total, as a percentage of GDP.
Positive and negative contributions stack away from zero independently, so the
visible column height equals the arithmetic sum of the layers.}
\label{fig:contributions}
\end{figure}

\subsection{Persistence, and why magnitudes are regime-specific}

Table~\ref{tab:composition} reports sign frequencies and magnitudes inside each
statistical regime.

\begin{table}[htbp]
\centering
\caption{Subsector composition of the general-government balance, by statistical regime. Magnitudes are reported inside each regime because the two differ enough that a pooled mean describes neither.}
\label{tab:composition}
\small
\fitwidth{%
\begin{tabular}{@{}llrrrrr@{}}
\toprule
Regime & Subsector & N & Positive & Negative & Mean (\% GDP) & Median (\% GDP) \\
\midrule
1977--1994 historical & General Government & 18 & 0 & 18 & -6.43 & -6.03 \\
1995--2025 modern & General Government & 31 & 4 & 27 & -4.04 & -4.17 \\
1977--1994 historical & Central Government & 18 & 0 & 18 & -6.61 & -6.31 \\
1995--2025 modern & Central Government & 31 & 0 & 31 & -4.67 & -4.33 \\
1977--1994 historical & Regional and Local & 18 & 5 & 13 & -0.27 & -0.19 \\
1995--2025 modern & Regional and Local & 31 & 13 & 18 & -0.09 & -0.10 \\
1977--1994 historical & Social Security Funds & 18 & 14 & 4 & 0.46 & 0.47 \\
1995--2025 modern & Social Security Funds & 31 & 29 & 2 & 0.72 & 0.61 \\
\bottomrule
\end{tabular}%
}
\end{table}


Three patterns are stable across both regimes. The Central Government balance is
negative in every observation of the panel, all \CentralNegativeYears\ of them. The
Social Security balance is positive in \SsfPositiveYears\ of \PanelYears, negative
in \SsfNegativeYears, with a longest unbroken positive run of
\SsfLongestPositiveRun\ years. Regional and Local Government is the only subsector
that changes sign with any regularity, positive in \RegionalPositiveYears\ years,
and it is small in both directions.

Magnitudes, by contrast, are strongly regime-dependent, and this is the reason the
table is split. The aggregate balance averages \GgMeanHistorical\% of GDP over
\PanelStart--1994 and \GgMeanModern\% over 1995--\PanelEnd. A pooled mean of
\GgMeanPooled\% falls between the two and corresponds to no period in the data.
Central Government moves from \CentralMeanHistorical\% to \CentralMeanModern\% and
the Social Security Funds from \SsfMeanHistorical\% to \SsfMeanModern\%.

Sign counts survive pooling far better than magnitudes, because a sign does not
depend on the level convention of a vintage. That asymmetry is why we pool the
counts in the discussion above and never pool the means.

\subsection{The years with a positive aggregate balance}

The aggregate balance is positive in \NumPositiveYears\ of \PanelYears\ years:
\PositiveYearList. Table~\ref{tab:positive} gives the composition of each.

\begin{table}[htbp]
\centering
\caption{Every year in which the aggregate balance is positive, with its subsector composition and the offset ratio.}
\label{tab:positive}
\small
\fitwidth{%
\begin{tabular}{@{}rrrrrrr@{}}
\toprule
Year & GG (M EUR) & Central (M EUR) & Regional/local (M EUR) & SSF (M EUR) & Non-SSF (M EUR) & Offset ratio \\
\midrule
2019 & 249 & -3,331 & 604 & 2,975 & -2,727 & 1.091 \\
2023 & 3,030 & -2,447 & -243 & 5,720 & -2,690 & 2.126 \\
2024 & 1,863 & -4,424 & 256 & 6,031 & -4,168 & 1.447 \\
2025 & 2,059 & -5,636 & 630 & 7,065 & -5,006 & 1.411 \\
\bottomrule
\end{tabular}%
}
\end{table}


The composition is the same in all \NumPositiveYears. The combined
non-Social-Security balance is negative, the Social Security balance is positive,
and the latter exceeds the former in absolute size --- which is precisely the
condition under which the offset ratio of Equation~\eqref{eq:offset} is defined and
exceeds one. In \LatestYear\ the identity reads
\[
\LatestGg = \LatestCentral + \LatestRegional + \LatestSsf \quad \text{M EUR},
\]
with a combined non-Social-Security balance of \LatestNonSsf\ million euro and an
offset ratio of \LatestOffset. In ratio terms the aggregate balance was
\LatestGgPct\% of GDP against a Social Security balance of \LatestSsfPct\%.

Figure~\ref{fig:offset} plots the ratio over the years in which it is defined.

\begin{figure}[htbp]
\centering
\includegraphics[width=\linewidth]{02_ssf_offset_ratio.png}
\caption{The offset ratio of Equation~\eqref{eq:offset}, for the years in which its
sign condition holds. The reference line at one marks the level at which the Social
Security balance equals the non-Social-Security shortfall. The ratio relates two
recorded quantities and is not a counterfactual.}
\label{fig:offset}
\end{figure}

We restate the caution of Section~\ref{sec:framework} because this is the result
most exposed to over-reading. That the ratio exceeds one in \LatestYear\ is a
statement about the relative size of two published numbers. It is not a statement
that the aggregate balance would have been negative under a different institutional
arrangement, and the identity contains nothing that could support such a statement.

\subsection{Annual movements}

Figures~\ref{fig:attributionmodern} and~\ref{fig:attributionhistorical} attribute
each annual change to the subsectors, on the GDP scaling of
Equation~\eqref{eq:scaled}. The two windows are drawn as separate panels and neither
crosses 1995: a single panel spanning the splice would place a vintage revision
among the economic movements and give it the same visual weight.

\begin{figure}[htbp]
\centering
\includegraphics[width=\linewidth]{03_balance_change_attribution.png}
\caption{Annual change in the General Government balance attributed to subsectors,
1996--\PanelEnd, as a percentage of current-year GDP. The black line is the identity
total.}
\label{fig:attributionmodern}
\end{figure}

\begin{figure}[htbp]
\centering
\includegraphics[width=\linewidth]{11_attribution_historical.png}
\caption{The same decomposition for 1978--1994, on the same scaling. The windows are
plotted separately because no change is computed through the 1995 splice.}
\label{fig:attributionhistorical}
\end{figure}

Table~\ref{tab:movements} ranks the \MovementsPerDirection\ largest improvements and
the \MovementsPerDirection\ largest deteriorations on that scaling, and attributes
each to a subsector and then to the revenue and expenditure movements behind it.

\begin{table}[htbp]
\centering
\caption{The largest annual improvements and deteriorations, ranked on the change scaled by current-year GDP, with the subsector and the revenue or expenditure movements that composed each one.}
\label{tab:movements}
\small
\fitwidth{%
\begin{tabular}{@{}lrrrrrrr@{}}
\toprule
Direction & Year & Change (\% GDP) & Central (M EUR) & Regional/local (M EUR) & SSF (M EUR) & From revenue (M EUR) & From expenditure (M EUR) \\
\midrule
Improvement & 2011 & 3.92 & 5,937 & 1,342 & -366 & 1,796 & -5,117 \\
Improvement & 2015 & 2.63 & 4,283 & 121 & 321 & 1,854 & -2,871 \\
Improvement & 2021 & 2.52 & 5,570 & -394 & 270 & 9,188 & 3,742 \\
Improvement & 1986 & 2.52 & 355 & 97 & 181 & 2,092 & 1,458 \\
Improvement & 2018 & 2.44 & 5,053 & 151 & -208 & 4,794 & -203 \\
Deterioration & 2009 & -5.98 & -8,939 & -289 & -1,268 & -3,802 & 6,694 \\
Deterioration & 2020 & -5.88 & -10,224 & -811 & -777 & -4,176 & 7,636 \\
Deterioration & 1981 & -5.40 & -385 & -48 & -45 & 573 & 1,051 \\
Deterioration & 1984 & -4.71 & -922 & -1 & 133 & 943 & 1,733 \\
Deterioration & 1980 & -4.02 & -233 & -34 & -29 & 523 & 819 \\
\bottomrule
\end{tabular}%
}
\par\smallskip
\begin{minipage}{\linewidth}\footnotesize\itshape 1995 is excluded: the 1994-to-1995 change straddles the vintage splice in both panels and mixes a statistical revision with an economic movement.\end{minipage}
\end{table}


The scaling changes what the ranking shows, which is the point of using it.
\HistoricalEpisodes\ of the ten episodes fall in the historical regime --- episodes
that a nominal ranking places nowhere near the top, because the euro amounts
involved are small relative to modern GDP. The largest single improvement is
\TopImprovementYear\ at \TopImprovementPct\ percentage points of GDP and the largest
deterioration is \TopDeteriorationYear\ at \TopDeteriorationPct.

Central Government accounts for the largest share of every one of the ten episodes.
Given that it also carries the level of the aggregate balance
(Section~\ref{sec:results}), this is less surprising than it first appears, and it
is worth stating what it does not mean: the subsector accounting arithmetically for
a movement has not been shown to have produced it.

\subsection{Social Security mechanisms}

We now turn to the second research question. Two distinct layers have to be kept
apart, and the value of the exercise lies in the separation.

\paragraph{The national-accounts layer.} The ESA 2010 Social Security Funds sector
is the entity that appears in Equation~\eqref{eq:identity}. In \LatestYear\ social
contributions were \ContributionSharePct\% of its total revenue.
Figure~\ref{fig:ssfshare} plots that share over the panel.

\begin{figure}[htbp]
\centering
\includegraphics[width=\linewidth]{05_ssf_contribution_share.png}
\caption{Social contributions as a share of total Social Security Funds revenue.
The line breaks at 1996--1999, where subsector components do not exist in either
source family.}
\label{fig:ssfshare}
\end{figure}

\paragraph{The budget-execution layer.} The CFP publishes the system decomposed into
the Previdential system, the Social Protection of Citizenship system and special
regimes. These are budget-execution aggregates on a different accounting boundary.
Figure~\ref{fig:ssfsystems} shows them.

\begin{figure}[htbp]
\centering
\includegraphics[width=\linewidth]{13_ssf_budget_systems.png}
\caption{CFP internal Social Security balances by system, in million euro. The
movement over the published window is concentrated in the Previdential system.}
\label{fig:ssfsystems}
\end{figure}

The movement is concentrated in one system. The Previdential balance goes from
\PrevidentialFirst\ million euro in \PrevidentialFirstYear\ to
\PrevidentialLatest\ million in \LatestYear, while the Citizenship balance stays
small in both directions, standing at \CitizenshipLatest\ million in \LatestYear.
Since contributions are the dominant revenue item of the Previdential system, and
since the contribution share of national-accounts revenue rose over the same years,
the two layers are telling a consistent story about where the movement sits. Neither
identifies why it happened, and we do not speculate: a contributory balance moves
with employment, wages, contribution rates, entitlement rules and demographics
simultaneously, and separating them needs a model this paper does not build.

\paragraph{Why the two layers must not be interchanged.} This is a methodological
result, and it is the one we would most want a reader to carry away.
Table~\ref{tab:ssfboundary} places the two side by side.

\begin{table}[htbp]
\centering
\caption{Two accounting boundaries for Social Security, reported side by side. The columns are never added, netted or reconciled.}
\label{tab:ssfboundary}
\small
\fitwidth{%
\begin{tabular}{@{}rrrrrr@{}}
\toprule
Year & ESA 2010 (M EUR) & Budget systems (M EUR) & Previdential (M EUR) & Citizenship (M EUR) & Difference (M EUR) \\
\midrule
2019 & 2,975 & 2,776 & 2,333 & 443 & 199 \\
2020 & 2,198 & 2,072 & 1,649 & 423 & 126 \\
2021 & 2,468 & 2,348 & 1,695 & 653 & 120 \\
2022 & 4,299 & 4,062 & 4,172 & -110 & 237 \\
2023 & 5,720 & 5,486 & 5,485 & 1 & 234 \\
2024 & 6,031 & 5,595 & 5,840 & -245 & 436 \\
2025 & 7,065 & 6,657 & 6,712 & -55 & 408 \\
\bottomrule
\end{tabular}%
}
\end{table}


In \LatestYear\ the ESA 2010 balance was \SsfEsaLatest\ million euro while the three
budget systems summed to \SsfBudgetLatest\ million, a difference of
\SsfBoundaryLatest\ million. Across the \SsfBoundaryYears\ overlapping years the
difference ranges from \SsfBoundaryMin\ to \SsfBoundaryMax\ million euro. It is
never zero.

Figure~\ref{fig:ssfboundary} makes the scale problem visible: on a shared axis the
two series are nearly identical, and the difference is only legible on an axis of
its own.

\begin{figure}[htbp]
\centering
\includegraphics[width=\linewidth]{14_ssf_accounting_boundary.png}
\caption{The two Social Security accounting boundaries. Upper panel: both balances
on a shared axis, where they are almost indistinguishable. Lower panel: the
difference on its own axis. The two are never added, netted or reconciled here.}
\label{fig:ssfboundary}
\end{figure}

The consequence is practical. A Social Security figure quoted from the budget
documents is not the figure that enters the national-accounts identity, and
substituting one for the other introduces an error of a few hundred million euro
without any signal that it has happened. We therefore report both and add neither.

\subsection{Primary balance: the headline sign is not the underlying sign}

The Central Government balance is negative in all \CentralNegativeYears\ years of
the panel. Read alone, that invites a strong conclusion: that the subsector runs an
underlying deficit throughout. The detailed accounts do not support it.

The primary balance is reconstructed as $PB_{i,t} = B_{i,t} + I_{i,t}$, where $I$ is
interest expenditure, and checked against the published primary-balance row.
Table~\ref{tab:primarysigns} contrasts the two sign frequencies.

\begin{table}[htbp]
\centering
\caption{Headline against primary balance sign frequencies, over the sector-years for which interest expenditure is published.}
\label{tab:primarysigns}
\small
\fitwidth{%
\begin{tabular}{@{}lrrrrr@{}}
\toprule
Subsector & N & Headline $<0$ & Primary $>0$ & Mean interest (\% GDP) & Peak interest (\% GDP) \\
\midrule
General Government & 49 & 45 & 22 & 4.16 & 7.89 \\
Central Government & 45 & 45 & 15 & 4.17 & 7.73 \\
Regional and Local & 45 & 29 & 20 & 0.13 & 0.62 \\
Social Security Funds & 45 & 6 & 39 & 0.00 & 0.04 \\
\bottomrule
\end{tabular}%
}
\end{table}


Over the \PrimaryObservedYears\ years for which interest expenditure is published,
the Central Government headline balance is negative in
\PrimaryHeadlineNegative\ --- all of them --- while the primary balance is positive
in \PrimaryPositiveYears: \PrimaryPositiveList. Those years are not confined to the
recent period; they run from \PrimaryPositiveFirst\ to \PrimaryPositiveLast. Interest
averaged \InterestMeanPct\% of GDP over the panel and peaked at
\InterestPeakPct\%, and that is the arithmetic separating the two series.
Figure~\ref{fig:primary} shows the primary balance crossing zero repeatedly while the
headline balance never does.

\begin{figure}[htbp]
\centering
\includegraphics[width=\linewidth]{06_central_primary_balance.png}
\caption{Central Government headline balance, primary balance and interest
expenditure, as shares of GDP. The primary balance crosses zero repeatedly; the
headline balance does not. The line breaks at 1996--1999.}
\label{fig:primary}
\end{figure}

Both statements are descriptive and they do not conflict. The correct summary is
that the Central Government B.9 is negative throughout the panel \emph{and} that its
primary balance is positive in a non-trivial minority of the observed years.

The symmetric caution applies. A positive primary balance is not a sustainability
result: it excludes interest by construction and therefore says nothing about the
debt path, which depends on the interest burden it removes. Neither series, on its
own, licenses a conclusion about sustainability.
